% Management in Global Markets - Week 14 Cornell Final Notes (Spring 2026, Mon/Tue)
% Compile: pdflatex management-global-markets-final-cornell-master-notes-mon-tue-spring2026.tex
\documentclass[10pt,a4paper]{article}
\usepackage[utf8]{inputenc}
\usepackage[T1]{fontenc}
\usepackage[margin=1.2cm]{geometry}
\usepackage{xcolor}
\usepackage{longtable}
\usepackage{enumitem}
\setlist[itemize]{leftmargin=1.2em,itemsep=1pt,topsep=2pt}
\definecolor{bizblue}{RGB}{20,52,110}
\definecolor{soft}{RGB}{245,248,252}
\setlength{\parindent}{0pt}
\setlength{\parskip}{0.35em}

\newcommand{\cornell}[3]{
\subsection*{#1}
\begin{longtable}{|p{0.22\linewidth}|p{0.74\linewidth}|}
\hline
\rowcolor{soft}\textbf{Cue / Questions} & \textbf{Notes} \\
\hline
#2 & #3 \\
\hline
\end{longtable}
}

\begin{document}
\begin{center}
{\Large\bfseries\color{bizblue} Management in Global Markets - Week 14 Final Cornell Master Notes}\\
{\small Spring 2026 (typical class days: Mon/Tue) | Consolidated from weekly topics, cases, and midterm revision threads}
\end{center}

\section*{Performance and Focus Rules}
\begin{itemize}
  \item \textbf{Business current marker:} C (5/10) = aprobado bajo (AP).
  \item \textbf{Priority rule:} this is a high-focus course for finals improvement.
  \item \textbf{Global fail thresholds:} WIP 3.49--0; UCM 4.99--0; suspenso = 4.
\end{itemize}

\section*{Cornell Notes by Exam Theme}

\cornell{1) Globalization, Slowbalization, and Strategic Context}
{What changed in global context?\\
How does slowbalization alter firm choices?\\
What matters for managers?}
{
\textbf{Global baseline:} firms coordinate across countries to access demand, capabilities, and efficiency.\\
\textbf{Slowbalization effect:} fragmentation, resilience concerns, and policy shifts reduce pure efficiency logic.\\
\textbf{Managerial implication:} strategy must balance growth with geopolitical and operational risk.\\
\textbf{Exam formula:} trend -> mechanism -> strategic adjustment.\\
\textbf{Key warning:} do not treat globalization as linear or irreversible.
}

\cornell{2) Internationalization Path and Entry Modes}
{Which market first?\\
Which mode and why?\\
How to evaluate risk-return?}
{
\textbf{Entry-mode spectrum:} exporting, licensing, partnerships, JV, FDI.\\
\textbf{Selection criteria:} control needs, resource commitment, institutional risk, speed.\\
\textbf{ICP/persona linkage:} international decisions should reflect segment-level demand fit.\\
\textbf{Strategic sentence:} ``Entry mode is a governance choice under uncertainty.''\\
\textbf{Typical mistake:} recommending one mode without capability and risk justification.
}

\cornell{3) Global Value Chains, Supply Chains, and Execution Risk}
{Where is value created?\\
What are bottlenecks?\\
How to improve resilience?}
{
\textbf{Value-chain view:} firms allocate activities across countries by capability-cost trade-offs.\\
\textbf{Supply-chain discipline:} visibility, supplier diversification, and process coordination reduce disruption cost.\\
\textbf{Risk categories:} political, currency, compliance, logistics, dependency concentration.\\
\textbf{Managerial task:} combine efficiency with continuity safeguards.\\
\textbf{Exam point:} resilience is not anti-efficiency; it is risk-adjusted performance design.
}

\cornell{4) Integration Areas, Mercosur, and Regional Strategy}
{Why regional blocs matter?\\
How do integration levels affect firms?\\
What changes in execution?}
{
\textbf{Regional integration logic:} trade barriers, standards, and movement rules shape market accessibility.\\
\textbf{Mercosur relevance:} illustrates institutional effects on expansion options and constraints.\\
\textbf{Strategic consequence:} region-specific legal/commercial setup affects GTM and operating model.\\
\textbf{Exam structure:} integration level -> firm impact -> recommended adjustment.\\
\textbf{Frequent error:} discussing regions abstractly without linking to firm decisions.
}

\cornell{5) Financial and Scenario Logic for Global Decisions}
{How to handle uncertainty?\\
What scenarios are realistic?\\
What levers are controllable?}
{
\textbf{Scenario method:} build plausible futures and pre-define response triggers.\\
\textbf{Financial levers:} pricing, cost base, channel mix, investment pacing, hedging exposure.\\
\textbf{Decision quality:} strongest answers show contingent plans, not one static recommendation.\\
\textbf{Exam line:} ``A robust international strategy preserves optionality under uncertainty.''\\
\textbf{Priority reminder:} spend extra finals time here due to current grade risk.
}

\cornell{6) EU Legal Order, Single Market, and Rule-of-Law Tools (Managerial Lens)}
{Why should global managers care about CJEU--national court cooperation?\\
What is a preliminary ruling in one sentence?\\
How do Article 7, infringement cases, and conditionality differ?\\
What is the ``double standard'' / inconsistency critique?\\
Why does ``constitutional identity'' matter for operating risk?\\
\textit{Briefing drill:} three sanctions---nature vs main weakness?}
{
\textbf{Identity tension (exam):} the EU is moving from a \textbf{primarily economic project} toward a \textbf{Union of Values} (Art.\ 2); the \textbf{inconsistency} / ``double standard'' critique flags an \textbf{identity crisis}: economic vs values enforcement can look \textbf{asymmetric} to capitals and publics.\\
\textbf{Bridge to global strategy:} same polity is a \textbf{single market} (uniform rules) and a stressed \textbf{political union}; shocks in one state spill across supply chains and procurement (\textbf{butterfly effect} in EU--business framings).\\
\textbf{Judicial ``high authority'' (metaphor):} the \textbf{CJEU} is not a national supreme court you only ``reach'' on appeal; it is the apex \textbf{interpretive} power whose doctrine binds through \textbf{preliminary rulings} and national judges' duty of cooperation.\\
\textbf{Preliminary ruling (Art.\ 267 TFEU):} a \textbf{horizontal bridge}---e.g.\ Madrid and Warsaw judges refer questions so both apply the \textbf{same} treaty meaning, curbing \textbf{legal fragmentation} of the single market (contrast a US-style vertical ladder to one federal top court).\\
\textbf{Primacy vs sovereignty:} EU law claims \textbf{supremacy} over conflicting national statutes (\textit{Costa} line); \textbf{illiberal} executives stress \textbf{constitutional identity}. \textbf{Juridification:} political fights become \textbf{technical court battles}, not only diplomacy---\textbf{compliance uncertainty} for investors.\\
\textbf{Rule-of-law toolbox (compare/contrast / essay):}\\[0.15em]
{\small\noindent\begin{tabular}{@{}p{0.26\linewidth}p{0.22\linewidth}p{0.44\linewidth}@{}}
\hline
\textbf{Mechanism} & \textbf{Nature} & \textbf{Key weakness} \\
\hline
Art.\ 7 TEU & Political (voting-rights suspension) & \textbf{Unanimity}; ally \textbf{loyalty pact} can block. \\
\hline
Infringement & Legal (Commission--CJEU) & Slow; fines may not force systemic political U-turn. \\
\hline
Conditionality & Economic (budget freeze) & Must tie to EU \textbf{financial interests}. \\
\hline
\end{tabular}}\\
\normalsize
\textbf{Inconsistency / ``double standard'' (e.g.\ Gros-type reading):} harder/faster \textbf{economic} enforcement than \textbf{values} (courts, media) can erode moral authority and fuel backlash---input to \textbf{country-risk} / reputational scoring.\\
\textbf{Operators and Art.\ 2:} due diligence on courts, enforcement, procurement---not tariffs alone.\\
\textbf{Exam one-liner:} ``Uniform CJEU doctrine protects scale; Art.\ 7 gridlock shifts leverage to \textbf{money} (conditionality), not votes.''
}

\cornell{7) Porter Value Chain, Hofstede, and IB/Human Resources (Lecture Thread)}
{What are Porter primary vs support activities?\\
Typical internationalization sequence across the chain?\\
Which Hofstede dimensions did the lecture stress?\\
Why did the French $\rightarrow$ Chile story fail?\\
\textit{Briefing:} cost fit \emph{vs} cultural fit?}
{
\textbf{Lecture anchor:} international \textbf{logistics} and relocation---professor's experience moving a \textbf{French} firm's footprint toward \textbf{Chile} and \textbf{Mexico}; globalization as \textbf{relocating whole departments}, not only shipping goods---match \textbf{cost} and \textbf{cultural fit}.\\
\textbf{Porter value chain (typology):} map internal activities to see where \textbf{value} and \textbf{margin} are created and where \textbf{cost} can fall; visualize interlocking activities (``cogs''), not isolated silos.\\
\textbf{Primary activities:} \textbf{inbound logistics} (inputs from suppliers) $\rightarrow$ \textbf{operations} (production) $\rightarrow$ \textbf{outbound logistics} (distribution) $\rightarrow$ \textbf{marketing \& sales} $\rightarrow$ \textbf{service} (after-sales).\\
\textbf{Support activities:} \textbf{procurement}, \textbf{human resources} (hire/train/retain), \textbf{technology development} (R\&D), \textbf{firm infrastructure} (finance, legal, general management).\\
\textbf{Internationalization sequence (lecture):} entering new markets often starts with \textbf{primary} front-end moves---especially \textbf{marketing \& sales}---then later \textbf{operations}/manufacturing may shift \textbf{offshore} if \textbf{factor costs} or access justify \textbf{FDI}.\\
\textbf{FDI reminder:} deploy \textbf{capital} abroad to control operations---bypass barriers, chase labor/cost advantages; still needs HR and culture fit.\\
\textbf{Hofstede ``human factor'':} culture can derail a perfect spreadsheet plan (risk of being ``fired''). \textbf{Power distance}---acceptance of hierarchy; lecture contrast \textbf{high} (e.g.\ Mexico, Spain) vs \textbf{low} (e.g.\ Nordic Europe, USA): challenge to authority, meeting style, delegation. \textbf{Individualism vs collectivism}---``I''/achievement vs ``We''/loyalty (lecture linked individualism to wealthier markets and collectivism to developing contexts---use as teaching shorthand, not stereotype in essays). \textbf{Masculinity vs femininity}---competition/achievement pole vs cooperation/quality-of-life pole.\\
\textbf{Class vocabulary:} \textbf{HR} owns employee lifecycle; \textbf{profitability} ties chain choices to financial outcomes; \textbf{hierarchy} reads differently under power distance; \textbf{tax havens} appear in offshore structuring debates; \textbf{Incoterms} (not ``Intercoms'') allocate buyer/seller risk in international sales contracts.\\
\textbf{Strategic warning (France $\rightarrow$ Chile):} \textbf{economic complexity} + \textbf{cultural ignorance} as failure traps---assuming a Paris-proven design ``exports'' without adapting to local \textbf{education}, \textbf{authority norms}, and feedback channels; business is \textbf{social science}, not only math.
}

\cornell{8) International Market Entry, B vs C, Culture Shock, and Hofstede Extensions (Lecture Thread)}
{What risk buckets apply when moving abroad?\\
How read GDP and political risk?\\
\textbf{B} vs \textbf{C} markets---which is often safer and why?\\
What is the honeymoon--shock arc?\\
Hofstede add-on: long- vs short-term orientation?}
{
\textbf{Global everyday hook:} \textbf{Vietnam}-made shoes, \textbf{China}-made shirts, \textbf{Botswana} gold---illustrate real \textbf{global supply chains} and the gamble when a firm ``moves'' beyond the home country.\\
\textbf{Market entry risk (taxonomy):} \textbf{cultural} (values, norms, religion); \textbf{legal/administrative} (tax codes, \textbf{trade barriers}, different legal systems); \textbf{economic} (\textbf{currency} swings, \textbf{GDP} levels and \textbf{growth}).\\
\textbf{Market indicators:} \textbf{GDP} as coarse \textbf{wealth/size} signal; \textbf{political risk} as stability/governance---avoid or discount countries with high \textbf{expropriation} or state-collapse fears.\\
\textbf{Trade barriers / protectionism:} tariffs/taxes shielding domestic industry; lecture note---high duties often hit goods the country \textbf{already produces} (e.g.\ agriculture).\\
\textbf{Competition strategy (B vs C):} \textbf{B} = many incumbents (``crowded''); \textbf{C} = no visible competitors. Counter-intuitive class point: \textbf{follow rivals into B} when it signals \textbf{proven demand} and existing \textbf{logistics}/ecosystem; \textbf{C} is often empty because \textbf{no demand} or \textbf{no infrastructure}---``safe'' absence of competition can be a trap.\\
\textbf{Rivals---operational triad to win share:} beat them on \textbf{price}, \textbf{quality}, or \textbf{delivery time} (pick where you can defend).\\
\textbf{Hofstede extensions:} \textbf{power distance} / \textbf{individualism--collectivism} reinforced (high PD in e.g.\ Mexico/Spain; \textbf{individualism} USA; \textbf{collectivism} linked to many Asian/African lecture examples---shorthand only). Add \textbf{long-term orientation} (future, saving, persistence---student ``mortgage/retirement'' motivation) vs \textbf{short-term} (present results). \textbf{Secularism} noted in \textbf{Turkey} historical context (state--religion separation).\\
\textbf{Culture shock model:} phase 1 \textbf{honeymoon/excitement}; phase 2 \textbf{shock} when rules, language, and management style diverge. \textbf{Professor tip:} you can optimize \textbf{taxes} and \textbf{logistics}, but not \textbf{communication} and \textbf{language}.\\
\textbf{Logistics geography:} \textbf{landlocked} states (e.g.\ Bolivia, Paraguay in South America examples) pay corridor/coordination premiums.\\
\textbf{Exam line:} pair \textbf{screening indicators} (GDP, political risk, competitive structure) with \textbf{culture} before celebrating a ``blue ocean'' that is merely \textbf{empty}.
}

\cornell{9) Cross-Cultural Management, PDI, CQ, and the U-Curve (Lecture Thread)}
{What does Hofstede \textbf{power distance} predict in the workplace?\\
\textbf{PDI} high vs low---behaviors and examples?\\
Three levels of \textbf{cultural learning}?\\
\textbf{93\% rule}---why it matters in negotiation?\\
Full \textbf{U-curve}---danger zone for expats?\\
\textit{Class activity:} self-assessment on hierarchy/authority.}
{
\textbf{Course blend:} \textbf{international business strategy} + \textbf{cross-cultural management}; culture gaps (especially \textbf{PDI}) drive ops, leadership style, and \textbf{expatriate} \textbf{culture shock}.\\
\textbf{Power distance (PDI):} degree to which less powerful members \textbf{accept} that power is \textbf{unequal}. \textbf{High PDI} (e.g.\ Mexico, Russia, Spain in lecture): steep \textbf{hierarchy}, young managers (``baby bosses'') over older staff, \textbf{micromanagement}, reluctance to contradict superiors. \textbf{Low PDI} (e.g.\ USA, Sweden, Iceland): flatter teams, accessible bosses, stronger \textbf{equality}/participation norms (lecture snapshot---indexes are coarse).\\
\textbf{Individualism vs collectivism (recap):} USA-style ``I''/achievement vs many Asian/African ``We''/harmony framings in lecture shorthand.\\
\textbf{Human factor / France $\rightarrow$ Chile:} failure framed as missing \textbf{cultural intelligence (CQ)}---management style from HQ did not \textbf{translate}; technical plan $\neq$ local legitimacy.\\
\textbf{Integration ladder (three levels):} (1) \textbf{Academic}---books/lectures (\textbf{low}); (2) \textbf{training/awareness}---cases, short mobility (\textbf{medium}); (3) \textbf{integration}---long immersion building \textbf{CQ} (\textbf{high}).\\
\textbf{Communication:} verbal vs \textbf{non-verbal}; lecture cites \textbf{65--93\%} of meaning from body language, tone, face---in negotiation, \textbf{how} something is said often beats \textbf{what} is said.\\
\textbf{U-curve (culture shock):} (1) \textbf{honeymoon} (excitement); (2) \textbf{shock/the low}---language, norms, isolation (\textbf{quit zone}); (3) \textbf{adjustment/intuition}---empathy and navigation; (4) \textbf{mastery/steady state}---comfort and effectiveness.\\
\textbf{Strategic insight---empathy:} arguable top global-manager skill is \textbf{empathy}, not credentials alone; in high-PDI settings read ``hidden rules'' of \textbf{chain of command}; ignoring undercurrents $\rightarrow$ very high failure odds in class framing.\\
\textbf{Bridges to other themes:} \textbf{sovereignty} reluctance toward IOs appears in sovereignty--integration debates; align with block 6--8 when writing \textbf{country + culture} risk together.\\
\textbf{Exam discipline:} pair Hofstede labels with \textbf{observable behaviors} (meetings, feedback, age/status) not only country names.
}

\cornell{10) Four Phases of Internationalization, Risk Lifecycle, and EU Mobility Context (Lecture Thread)}
{What are the \textbf{four phases} from conception to establishment?\\
Why skip phases 1--2 before \textbf{execution}?\\
\textbf{Artificial} vs \textbf{natural} risk---what can be ``escaped''?\\
Where does \textbf{risk peak} on the lifecycle graph?\\
\textit{Timeline:} \textbf{six years} to maturity---why do ``middle'' firms fail?\\
Schengen, rule of law, backsliding---how mentioned in discussion?}
{
\textbf{Phase overview (internationalization):} (1) \textbf{Conception}---internal capacity and fit; \textbf{SWOT}-style scan of strengths/weaknesses vs external opportunity. (2) \textbf{Market study}---deep read of target environment; lecture: skipping (1)--(2) invites \textbf{catastrophic failure} in phase 3. (3) \textbf{Execution}---physical move: plants, logistics, people; about \textbf{two years}; apex \textbf{financial} and operational exposure (cash out before revenues stabilize). (4) \textbf{Establishment}---\textbf{steady state}: settled operations, profitability emerging, \textbf{risk} trajectory finally slopes down.\\
\textbf{Risk lifecycle diagram:} reward--risk narrative peaks in the transition from \textbf{planning/study} into \textbf{execution}---the ``valley of death'' for under-capitalized or impatient sponsors.\\
\textbf{Artificial vs natural risk:} \textbf{Artificial}---systems you can engineer around or partially ``escape'' (\textbf{logistics}, legal/tax rules, administrative friction). \textbf{Natural}---\textbf{language}, \textbf{communication}, culture: nearly impossible to bypass without long \textbf{integration}/\textbf{CQ}; echoes block 8--9 (tax optimization $\neq$ cultural fix).\\
\textbf{FDI in execution:} acquisitions, \textbf{JV}, greenfield---capital commitment lands hardest in phase 3.\\
\textbf{``Middle business'' / patience thesis:} full \textbf{maturity} often $\approx$ \textbf{six years}; many failures reflect inability to fund/hold through the \textbf{two-year execution} trench when outflows dominate inflows.\\
\textbf{Discussion bridges (EU):} \textbf{Schengen Area} ($\sim$29 states, internal open borders) illustrates integration benefits for mobility (e.g.\ student travel Madrid--Paris). Conversation tied \textbf{democratic backsliding} and \textbf{rule of law} (e.g.\ Hungary/Poland) to wider integration debates; \textbf{incentives} (e.g.\ EU support for newer members) as policy levers. \textbf{Excessive bureaucracy} cited as a Southern Europe business friction.\\
\textbf{Study-abroad ops aside (2026):} financial \textbf{literacy} under stress---housing \textbf{security deposit} refunds, university \textbf{payroll/bursar} routing, remote-work POS tools, \textbf{Spain vs US} cost-of-living comparison---parallel to how \textbf{artificial} frictions differ from \textbf{natural} adaptation needs.\\
\textbf{Migration literacy:} \textbf{Fujian}/\textbf{Fuzhou} (China) referenced as personal/global ``small world'' identity link in class banter---optional color, not core exam definition.\\
\textbf{Exam takeaway:} sequence phases before recommending \textbf{speed}; map risks to artificial (manage/hedge) vs natural (people/time) buckets.
}

\cornell{11) International Project Management, HR Staffing, and Margin Control (Lecture Thread)}
{\textbf{Expat} vs \textbf{local}---when HQ staff, when locals?\\
Why \textbf{commitment} over raw credentials?\\
How track a USD \textbf{\$5M} project from kickoff to close?\\
\textbf{Green} vs \textbf{red} projects on the margin dashboard?\\
Operational risks in developing markets---beyond the spreadsheet?\\
\textit{Trust:} who gets ``full control'' of capital abroad?}
{
\textbf{Lecture anchor:} professor's \textbf{industrial} experience in \textbf{Mexico}---ran a division of $\sim$\textbf{3,000} employees; illustrates \textbf{financial} and \textbf{human} complexity of internationalization at operating scale.\\
\textbf{Expatriates vs local hiring (strategy):} \textbf{HQ}/expat staff often justified for \textbf{high-tech} or sensitive programs---protect parent interests, enforce \textbf{technology transfer}, embed trusted know-how and \textbf{financial} controls; downside: \textbf{expats are expensive}. Locals add \textbf{lower cost} and \textbf{market/context} knowledge---blend typical.\\
\textbf{Commitment rule:} lecture prioritizes \textbf{commitment} over mere \textbf{education/experience}; a young, highly committed worker can beat an experienced but disengaged manager for \textbf{profitability} trajectory.\\
\textbf{Project financials (margins):} monitor \textbf{from start through completion}; example project \textbf{\$5M USD}; \textbf{20\%} margin respectable at kickoff; strong cost/ops discipline can stretch toward \textbf{26\%}. Tie reporting to \textbf{income statement} and \textbf{balance sheet} hygiene---especially moving \textbf{phase 3 (execution)} $\rightarrow$ \textbf{phase 4 (establishment)} (see block 10).\\
\textbf{Green vs red projects:} \textbf{green}---margins \textbf{held or improved} (e.g.\ 15\% stable or 20\% $\rightarrow$ 26\%); \textbf{red}---margin \textbf{compression} triggers review of \textbf{supply chain}, \textbf{logistics}, \textbf{labor} inputs. \textbf{Portfolio view:} healthy \textbf{average margin} (lecture example \textbf{19\%}) can still signal success even if single projects wobble.\\
\textbf{Operational risks (field):} \textbf{physical/environment} friction (lecture color: harsh sites, e.g.\ desert hazards delaying work); \textbf{safety/human} costs if training/supervision slip; \textbf{corruption}---managerial \textbf{integrity} framed as core ``asset'' when handling large \textbf{cash} and vendor flows.\\
\textbf{Strategic closing---trust \& control:} grant \textbf{full control} only to leaders with \textbf{implicit trust}; technical tasks can be \textbf{sourced locally}, but \textbf{loyalty to HQ} and \textbf{capital protection} rank first for international project leads in this framing.\\
\textbf{Exam link:} pair staffing choices with \textbf{phase}, \textbf{margin} color, and \textbf{risk register} (financial + people + integrity).
}

\cornell{12) CRM, Cash Flow, Down Payments, and Margin Portfolio Colors (Lecture Thread)}
{How does the firm \textbf{generate cash}?\\
\textbf{Down payment} rule before \textbf{execution}---why?\\
Portfolio: \textbf{Mitsubishi}, ``\textbf{Big One},'' \textbf{red} vs \textbf{yellow}?\\
\textbf{Customer retention} / CRM link to margins?\\
\textbf{6-month} reality vs long-range plans in Mexico/Chile?\\
\textit{Sovereign agent:} what are you hired to protect?}
{
\textbf{Thread:} \textbf{international business management} + \textbf{sales}/account strategy; Mexico industrial \textbf{spreadsheet} of projects as evaluation lens (extends block 11 numbers with \textbf{CRM} and \textbf{cash} discipline).\\
\textbf{Cash generation (two levers):} (1) \textbf{minimize expenses}/cost; (2) \textbf{maximize sales}/revenue---together feed \textbf{cash} the firm needs to survive \textbf{cycles}.\\
\textbf{``Cash is King'' / down payment:} large \textbf{industrial} jobs should not fully \textbf{execute} production without upfront \textbf{cash}; lecture cites about \textbf{50\%} \textbf{down payment} typical---starting ``naked'' exposes \textbf{capital} if the buyer stalls or defaults (\textbf{insolvency} risk for vendor).\\
\textbf{Margins \& accountability:} \textbf{margin} $=$ \textbf{revenue} $-$ \textbf{cost}; protect target---\textbf{excellence} lifts e.g.\ \textbf{20\%} $\rightarrow$ \textbf{26\%}; case row \textbf{\$60M USD} project at \textbf{20\%} margin (``Big One'') framed as \textbf{survival}/\textbf{cash-flow} anchor (``food and burgers'' for the company); stable \textbf{Mitsubishi} line at \textbf{20\%} as moderate baseline account.\\
\textbf{Traffic-light spreadsheet (lecture):} \textbf{red} row e.g.\ \textbf{11\%} margin---still above zero but \textbf{tight}; watch \textbf{labor} and \textbf{materials}. \textbf{yellow}---manager ``\textbf{lost the money},'' margin \textbf{15\%} or \textbf{below}; triggers \textbf{HQ} scrutiny and \textbf{termination} risk (color thresholds are course examples---map to your case numbers).\\
\textbf{CRM / retention:} one-off wins insufficient; \textbf{``coming back''} customers cut \textbf{acquisition} cost and stabilize the \textbf{portfolio}; ties \textbf{project execution} quality to future \textbf{revenue}.\\
\textbf{Commitment (echo):} \textbf{loyalty}/\textbf{hunger} beats lazy experience for high-stakes \textbf{international} seats (see block 11).\\
\textbf{Rule-of-law / integrity asset:} refusal of \textbf{bribes}, safeguarding \textbf{HQ capital} in \textbf{high-corruption} contexts builds rare \textbf{trust}.\\
\textbf{Short-term uncertainty:} strategy may aim \textbf{long term}, but operations often run in $\sim$\textbf{6-month} windows---\textbf{politics} and \textbf{regulation} in e.g.\ \textbf{Mexico}/\textbf{Chile} can freeze \textbf{investment} overnight; pair with \textbf{scenario} triggers (block 5).\\
\textbf{Sovereign-agent framing:} managers act as \textbf{agents} of \textbf{HQ owners}---job is \textbf{capital protection} and \textbf{return}, not sentiment; losing money ends careers regardless of ``kindness'' in lecture moral.\\
\textbf{Exam link:} connect \textbf{cash} terms, \textbf{milestone billing}/\textbf{down payment}, \textbf{margin band}, and \textbf{account} health in one answer.
}

\cornell{13) Economic Complexity, Location Choice, and the Firm’s Integration Matrix (Lecture Thread)}
{What does \textbf{ECI} / \textbf{economic complexity} claim about wealth?\\
\textbf{Stuck in the middle}---what does it take to escape?\\
\textbf{Diversification} vs \textbf{specialization} at country and firm level?\\
\textbf{Exporting} vs \textbf{internationalization}---risk ladder?\\
Why is \textbf{where} often harder than \textbf{how}?\\
\textbf{Taiwan}/\textbf{China}---why complexity raises \textbf{geopolitics}?}
{
\textbf{Core claim:} power and prosperity track not only \textbf{export volume} but \textbf{product complexity} and \textbf{diversity} of know-how embedded in what a country sells.\\
\textbf{ECI-style index logic:} rank economies by \textbf{knowledge intensity} of their export basket---complex goods (e.g.\ imaging equipment, semiconductors) vs \textbf{commodities} (e.g.\ bananas, bulk \textbf{oil}); complex producers tend to be more \textbf{resilient} and influential than one-good stories.\\
\textbf{``Stuck in the middle'' trap:} shallow, simple export mixes swing with prices; \textbf{moving up} needs \textbf{technology}, \textbf{skilled human capital}, and \textbf{stable institutions}/logistics (lecture bundle).\\
\textbf{North--South divide (lecture map):} most \textbf{high-value} global business clusters \textbf{north} of a stylized line; low-complexity / commodity \textbf{dependence} correlates with weaker infrastructure and narrower capabilities (\textbf{landlocked} states often pay extra \textbf{logistics} tax---see vocabulary).\\
\textbf{Diversification vs specialization:} \textbf{diversification} spreads a \textbf{wide basket} to reduce cyclical risk; \textbf{specialization} bets on \textbf{high-sophistication} niches (lecture cites \textbf{Switzerland} pharmaceuticals/watches, \textbf{Taiwan} semiconductors). Strategy choice, not one-size-fits-all.\\
\textbf{Exporting vs internationalization:} \textbf{exporting}---ship finished goods across borders (typically \textbf{lower} committed risk); \textbf{internationalization}---relocating \textbf{departments}/plants/supply links (\textbf{higher} capital, HR, and institutional demands).\\
\textbf{Location choice $>$ mode theater:} in complex markets, \textbf{fit} of host \textbf{competencies} to your \textbf{product sophistication} dominates picking ``the cheapest port''; local \textbf{knowledge cluster} must exist (or be built slowly)---you cannot parachute a frontier plant into a \textbf{competency void}.\\
\textbf{Firm matrix (lecture):} \textbf{vertical integration}---control sequential stages A$_1\ldots$A$_n$; \textbf{horizontal expansion}---place discrete \textbf{products} (e.g.\ product~8) in the \textbf{best location} for skills/demand (e.g.\ Switzerland for precision lines).\\
\textbf{Geopolitics link:} Taiwan/China salience framed by \textbf{advanced} products and supply choke points, not population \textbf{size} alone---if they exported only \textbf{fruit}, strategic stakes would read lower in this narrative.\\
\textbf{Exam discipline:} pair \textbf{country ECI intuition} + \textbf{firm matrix} + \textbf{market selection} (block 8/10) instead of GDP-only screening.
}

\cornell{14) Strategy Levels, Corporate Governance, and the ``Realist'' Manager (Lecture Thread)}
{\textbf{Corporate} vs \textbf{competitive} vs \textbf{functional} strategy---what question each answers?\\
\textbf{Realist triad:} executive power, financial control, emotional resilience?\\
How does \textbf{SWOT} sit before international moves?\\
\textbf{IP} risk when technology moves abroad?\\
\textit{Ambition trap:} vision \emph{vs} multi-year \textbf{price of the mission}?}
{
\textbf{Realist framing:} \textbf{internationalization} is a \textbf{high-risk}, \textbf{long-horizon} commitment---success needs a blend of \textbf{executive power} (authority to decide), \textbf{financial control} (guard \textbf{capital}, budgets, \textbf{margins}), and \textbf{emotional resilience} (withstand shock, isolation, and slow payoffs).\\
\textbf{Three strategy levels (pyramid):} \textbf{corporate}---which \textbf{industries} and \textbf{scope} (what businesses to be in); \textbf{competitive (business)}---how to \textbf{win} in a market (\textbf{price}, \textbf{quality}, differentiation---echo block 8 triad); \textbf{functional}---how \textbf{HR}, \textbf{finance}, \textbf{marketing}, ops \textbf{support} the chosen strategy. Aim: \textbf{synergy}---portfolio/business unit combinations worth more \textbf{together} than separately (lecture definition).\\
\textbf{Staffing dilemma (echo):} \textbf{HQ}/\textbf{expat} oversight protects \textbf{technology} and \textbf{capital}; \textbf{local} depth adds cost and context; expensive expats buy \textbf{governance}.\\
\textbf{Commitment beyond IQ:} baseline education not enough---\textbf{motivation} and \textbf{resilience} decide who lasts through \textbf{culture shock} and \textbf{years} overseas; ``\textbf{ambition trap}'' = polished \textbf{vision talk} without paying the \textbf{10-year} social/emotional \textbf{price of the mission}.\\
\textbf{Technology \& IP:} offshore moves raise \textbf{IP theft}/\textbf{copying} if legal controls and oversight lapse; lecture failure story: locals \textbf{replicated} know-how when guardrails failed.\\
\textbf{Financial guardianship:} international lead as \textbf{capital steward}---keep \textbf{income statements} healthy and turn \textbf{phase 3 execution} spend into \textbf{phase 4 establishment} growth (block 10 link).\\
\textbf{SWOT audit:} before scaling abroad, map \textbf{internal} \textbf{S}/\textbf{W} (e.g.\ patent vs no intl experience) and \textbf{external} \textbf{O}/\textbf{T} (e.g.\ \textbf{Chile} opportunity vs aggressive incumbents); aligns with \textbf{conception} phase (block 10).\\
\textbf{Closing aphorism (lecture):} ``\textbf{Cash is King},'' but \textbf{execution is the kingdom}---capital without disciplined delivery does not compound.\\
\textbf{Exam link:} stack \textbf{level of analysis} + \textbf{SWOT} + \textbf{people/IP/cash} controls in one coherent memo.
}

\cornell{15) Internationalization Spectrum, FDI Motives, and Export Mechanics (Lecture Thread)}
{\textbf{Internationalization} vs \textbf{exporting}---what counts as ``true'' int'l?\\
\textbf{Triple scope} alignment---corporate, competitive, functional?\\
\textbf{Market} vs \textbf{resource} vs \textbf{efficiency} seeking---examples?\\
Spectrum: \textbf{domestic} $\rightarrow$ \textbf{export} $\rightarrow$ \textbf{branch}/subsidiary $\rightarrow$ \textbf{global integration}?\\
\textbf{Direct} vs \textbf{indirect} export?\\
\textit{Shield:} how many years of burn can capital cover?}
{
\textbf{Definition (lecture):} \textbf{internationalization} is not merely \textbf{exporting}; it is \textbf{physically relocating} a \textbf{department} or \textbf{facility} (\textbf{marketing}, \textbf{production}, \textbf{HR})---needs \textbf{people}, \textbf{capital}, and host \textbf{infrastructure}.\\
\textbf{Exporting vs establishing abroad:} \textbf{exporting}---sell from \textbf{home base}; generally \textbf{lower} risk and \textbf{lower} operational control. \textbf{Branch}/\textbf{subsidiary}---legal and physical \textbf{presence}; \textbf{higher} cost/risk but deeper \textbf{penetration}/control (overlaps block 2--3, 13).\\
\textbf{Triple-scope check (consistency):} same three levels as block 14---\textbf{corporate} (\textbf{where} to go / vision), \textbf{competitive} (beat locals---\textbf{cheaper}? \textbf{better quality}?), \textbf{functional} (how \textbf{HR}/\textbf{finance} execute); misalignment breaks execution.\\
\textbf{Why internationalize (three motives):} \textbf{market seeking}---escape \textbf{saturated}/\textbf{collapsing} home demand, find customers; \textbf{resource seeking}---inputs, \textbf{cheap labor}, or \textbf{talent}; \textbf{efficiency seeking}---tax/\textbf{logistics}/cost basins that improve margin.\\
\textbf{``Stuck in the middle'' (echo):} expansion without \textbf{money} or \textbf{distinctive competencies} fails; need a \textbf{competitive advantage} locals cannot \textbf{quickly imitate} (block 13).\\
\textbf{Location choice (again):} poorest country pick can \textbf{wipe capital}---\textbf{French} firm in \textbf{Chile} cautionary tale (blocks 7, 9, 10).\\
\textbf{Spectrum of distance/complexity:} (1) \textbf{domestic}---safe, limited growth; (2) \textbf{exporting}---first cross-border step; (3) \textbf{branch}/\textbf{subsidiary}---moved departments; (4) \textbf{global integration}---coordinate many countries toward \textbf{synergy} (block 14).\\
\textbf{Export channels:} \textbf{direct export}---sell to \textbf{end client}; \textbf{indirect export}---via \textbf{middlemen}/distributors (risk, margin, and control trade-offs).\\
\textbf{Commitment bundle:} \textbf{long-term game}---\textbf{CQ} so a \textbf{France}-ready plan is adapted for \textbf{Mexico}/\textbf{Russia}; \textbf{will} to stay $\sim$\textbf{10+} years; \textbf{financial shield} to fund \textbf{2--3} early years before revenue fully carries operations (bridges blocks 9--12).\\
\textbf{Tax note (lecture):} host/home \textbf{income tax} can dominate location math---cited \textbf{French}-scale burdens (illustrative; verify in case numbers).\\
\textbf{Exam link:} sequence \textbf{motive} $\rightarrow$ \textbf{mode}/spectrum position $\rightarrow$ \textbf{location} screen $\rightarrow$ \textbf{cash} runway.
}

\cornell{16) IHRM, Motivation Theory, and Global vs Multi-domestic HR (Lecture Thread)}
{\textbf{Global} vs \textbf{multi-domestic} strategy---HR implications?\\
\textbf{Commitment} vs raw competency in staffing abroad?\\
\textbf{Maslow} and \textbf{Vroom}---what does each add for \textbf{expats}?\\
\textbf{Diverse} international teams---trade-offs?\\
\textbf{Executive power} / \textbf{negotiator authority}---why deals fail?\\
\textit{Ambition:} salary in a ``\textbf{disaster zone}'' without \textbf{CQ}?}
{
\textbf{Strategy--HR bridge:} \textbf{global} approach---high \textbf{standardization}; same product/\textbf{marketing} footprint (lecture examples: \textbf{iPhone}, \textbf{Coca-Cola}) with tight \textbf{corporate} norms. \textbf{Multi-domestic}---high \textbf{adaptation}; different offers, prices, \textbf{channels} per country (consumer goods, local services); needs deeper \textbf{local} HR depth and more \textbf{fragmented} people processes (ties \textbf{glocalization}/\textbf{trade-off} vocabulary).\\
\textbf{Staffing thesis (echo block 11):} \textbf{HQ} often weights \textbf{commitment} and \textbf{trust} over pure technical \textbf{pedigree}; \textbf{expatriates} (\textbf{``HQ'' alignment}) defend \textbf{capital} and \textbf{know-how}---cheaper to \textbf{move} a loyal lead than to \textbf{lose control} to misaligned locals in this framing (blend still typical; lecture is cautionary).\\
\textbf{Price of the mission (echo blocks 9--10, 14):} \textbf{internationalization} as $\sim$\textbf{10+}-year \textbf{game}; success needs \textbf{resilience} through \textbf{culture shock} and \textbf{CQ}; many exit from weak \textbf{motivation}, not lack of slides.\\
\textbf{Maslow (1943):} \textbf{hierarchy} from survival needs toward \textbf{esteem}/\textbf{recognition} and \textbf{self-actualization}; IHRM angle: for \textbf{expats}, upper tiers (recognition, growth, purpose) often rival \textbf{base pay} in retention---salary alone does not substitute.\\
\textbf{Taylor (scientific management)---backdrop:} classical \textbf{task} and \textbf{time} engineering; lecture positions \textbf{people-centered} motivation lenses (\textbf{Maslow}, \textbf{Vroom}) as what firms use today for \textbf{talent attraction}, not stopwatch-only shop-floor logic.\\
\textbf{Vroom (1964) expectancy theory:} three judgments before accepting/performing a hard assignment---(1) \textbf{expectancy}: effort $\rightarrow$ performance possible here? (2) \textbf{instrumentality}: if I perform, will \textbf{HQ} actually deliver promised \textbf{rewards}? (3) \textbf{valence}: do I \textbf{want} that reward (\textbf{money}, \textbf{promotion}, lifestyle)? Weak links demotivate faster abroad.\\
\textbf{International team diversity:} \textbf{pros}---creativity, \textbf{problem-solving}, access to varied \textbf{markets}/\textbf{norms}; \textbf{cons}---higher \textbf{turnover}, \textbf{communication} drag, internal \textbf{friction} across national backgrounds---needs explicit \textbf{team design} and norms.\\
\textbf{Executive authority / deal legitimacy:} \textbf{Mexico}/\textbf{Chile} ``war stories''---local counterparts without \textbf{respected} \textbf{authority}/\textbf{status} fail to \textbf{close}; \textbf{negotiation} outcomes depend on who credibly \textbf{signs} and who the counterparty \textbf{trusts} (links \textbf{power distance}, \textbf{trust}, block~14 \textbf{executive power}).\\
\textbf{Closing critique (lecture):} many claim \textbf{global ambition} but balk at the \textbf{personal price}: prolonged absence from home, \textbf{stress}, and \textbf{integrity} loads in tough environments; young \textbf{high salary} (e.g.\ illustrative \textbf{60k} at \textbf{27}) may \textbf{attract} but cannot replace \textbf{CQ} if the market is a \textbf{disaster zone}.\\
\textbf{Exam link:} pair \textbf{strategy archetype} $\rightarrow$ \textbf{staffing} $\rightarrow$ \textbf{motivation mechanism} (\textbf{Maslow}/\textbf{Vroom}) $\rightarrow$ \textbf{risk} (\textbf{diversity}/\textbf{authority}).
}

\cornell{17) Applied Studio: Strategic Business Plan---Caf\'e Ol\'e (Madrid) (Student Session)}
{\textbf{Pivot:} from professional \textbf{co-working} to \textbf{student}-centric hub?\\
\textbf{SWOT} + \textbf{PESTEL}---what did the team stress?\\
\textbf{Pricing} and \textbf{tiered membership}---trade-offs?\\
\textbf{``Slow is fast:''} \textbf{15 years} not \textbf{15 days}?\\
\textbf{Q1--Q4} rollout---what belongs in each quarter?\\
\textit{Exam craft:} \textbf{verbatim} instructor phrasing?}
{
\textbf{Source note:} collaborative \textbf{student} draft applying management/strategy frameworks to a real \textbf{Chamber\'i} venue---use as \textbf{pattern} for written plans, not factual claims about the operating business unless verified.\\
\textbf{Objective:} reposition \textbf{Caf\'e Ol\'e} around \textbf{Erasmus}/study-abroad students seeking \textbf{community} and reliable \textbf{study} space versus generic professionals alone.\\
\textbf{Pricing thread:} cited \textbf{current} pay-per-use roughly \textbf{\$5}/hour or \textbf{\$25}/day; \textbf{student}-oriented options included \textbf{20\%} hourly discount and about \textbf{EUR~30}/month \textbf{membership} (numbers are draft-level---recheck for feasibility).\\
\textbf{Competitive story:} many Madrid \textbf{caf\'es} impose \textbf{laptop bans} at peak or offer weak \textbf{Wi-Fi}; draft \textbf{differentiation} = guaranteed \textbf{workspace}, \textbf{quiet zones}, \textbf{printers}, amenity bundle (contrast \textbf{libraries}/crowding and private offices).\\
\textbf{SWOT (student synthesis):} \textbf{S}---fast internet, quiet zones, printers; \textbf{W}---incumbent \textbf{professional} positioning may \textbf{alienate} students; \textbf{O}---\textbf{university} partnerships, \textbf{language} exchanges, events; \textbf{T}---\textbf{coffee}/input inflation, \textbf{rent}, rival \textbf{response}.\\
\textbf{PESTEL (Spain/Madrid sketch):} \textbf{P}---EU labor rules, \textbf{student visa}/work-study policy; \textbf{E}---Eurozone \textbf{inflation} pressuring snack/\textbf{coffee} costs; \textbf{S}---\textbf{Erasmus} culture, demand for belonging; \textbf{T}---\textbf{fiber}/\textbf{5G} reliability for remote study; \textbf{En}---plastic/waste tension under ``\textbf{unlimited snacks}'' style perks; \textbf{L}---\textbf{zoning}/use rules for \textbf{business caf\'es} in Madrid.\\
\textbf{Contingency:} class-style \textbf{Plan B} examples---backup \textbf{internet} (\textbf{fiber} + \textbf{hotspot}), hedges for \textbf{utility} cost spikes; aligns with \textbf{scenario}/mitigation discipline (block 5, vocabulary).\\
\textbf{Executive summary spine:} \textbf{tiered membership} + \textbf{community events} bridges gap between costly private offices and \textbf{public libraries}; draft claimed ~\textbf{15\%} \textbf{loyalty} lift and filling \textbf{slow} morning slots via targeted \textbf{marketing} + university ties (verify in modeling).\\
\textbf{Rollout calendar (year one, draft):} \textbf{Q1}---\textbf{market research}, university \textbf{outreach}/partnerships; \textbf{Q2}---launch \textbf{student tier}, \textbf{social media}; \textbf{Q3}---\textbf{peak-hour} data, \textbf{quiet-zone} tweaks; \textbf{Q4}---\textbf{annual budget} review, expand \textbf{events}.\\
\textbf{Strategy aphorism (verbatim exam tip):} instructor line echoed in session---design strategy for \textbf{15 years}, not \textbf{15 days}; consistency beats one-off \textbf{sprints}.\\
\textbf{Meta---study behavior:} students cross-checked \textbf{audio} for exact professor \textbf{terminology}; signals \textbf{grade-sensitive} alignment between \textbf{artifacts} (\textbf{SWOT}/\textbf{PESTEL}/summary) and \textbf{rubric language}.\\
\textbf{Sidebar:} bursar/\textbf{deposit} logistics (e.g.\ \textbf{\$2{,}550}) cited in adjacent chat as parallel ``student admin'' friction---optional tie to \textbf{international student} cash-flow realism (block 10 EU ops aside).\\
\textbf{Exam link:} show you can \textbf{populate} matrices (not just name them), tie \textbf{PESTEL rows} to \textbf{pricing}/\textbf{ops} decisions, and state \textbf{assumptions} explicitly.
}

\cornell{18) Global vs Multi-domestic vs Transnational; Integration Ladder; ECI; Geopolitical Map (Lecture Thread)}
{\textbf{Global} / \textbf{multi-domestic} / \textbf{transnational}---how do they differ?\\
\textbf{Integration} stages from \textbf{FTA} to \textbf{EMU}?\\
\textbf{ECI}---why cite \textbf{Switzerland}?\\
\textbf{FDI} location---cheap labor alone?\\
\textbf{Ukraine} and EU---common exam barriers?\\
\textbf{Triad}, \textbf{G20}, north--south ``line''?\\
\textit{Closing:} \textbf{motivation} vs tech/capital; asking questions?}
{
\textbf{Strategy archetypes (expanded):} \textbf{global}---high \textbf{standardization} (\textbf{Germany}-style/industrial brands in shorthand) with ``think global, act local'' pockets; \textbf{multi-domestic}---deep per-country \textbf{adaptation} (\textbf{prices}, \textbf{marketing}); \textbf{transnational}---\textbf{hybrid} aiming for \textbf{both} high \textbf{global coordination/integration} and high \textbf{local responsiveness} (\textbf{complex}, costly to run---links \textbf{global--local trade-off}, block~16).\\
\textbf{Economic complexity (ECI echo):} index-style read of \textbf{knowledge intensity} in exports; \textbf{Switzerland} highlighted for \textbf{high-value}, \textbf{high-tech} baskets (\textbf{medical technology}, specialized \textbf{engineering})---prosperity tied to \textbf{sophistication}, not volume alone (block~13).\\
\textbf{Integration ladder (lecture stages):} (1) \textbf{free trade area (FTA)}---members drop \textbf{internal tariffs} while keeping own external policy (e.g.\ \textbf{USMCA}/NAFTA); (2) \textbf{customs union}---\textbf{common external tariff}; (3) \textbf{common market}---adds free movement of \textbf{labor} and \textbf{capital}; (4) \textbf{economic \& monetary union}---shared \textbf{currency} (\textbf{euro} zone) plus deeper macro rules---\textbf{EU} as peak reference in class narrative.\\
\textbf{Integration effects:} \textbf{static}---immediate trade pattern shifts (\textbf{trade creation} vs \textbf{trade diversion}); \textbf{dynamic}---longer-run \textbf{competition}, \textbf{scale}, productivity gains from bigger integrated markets.\\
\textbf{FDI \& location:} firms seek \textbf{competencies}---skill bases that can produce \textbf{complex goods}, not only the \textbf{cheapest} unskilled labor; aligns \textbf{competitiveness} with \textbf{human capital} depth (block~13 competency fit).\\
\textbf{IHRM bridge:} ``\textbf{internationalized people}'' prerequisite---\textbf{talent management}, specialist skills, and sustained \textbf{motivation}; echoes blocks~11, 16 that \textbf{people} carry the plan.\\
\textbf{Geopolitical map:} recurring north--south \textbf{high-value flow} metaphor (\textbf{95\%} band in lecture shorthand---illustrative); \textbf{triad} historically \textbf{US}/\textbf{EU}/\textbf{Japan}, now \textbf{shifting} weight toward \textbf{China}/\textbf{India}; \textbf{G20} cited for $\sim$\textbf{75\%} of \textbf{population} and bulk of \textbf{global trade}---use as coarse \textbf{power} context, not fine forecasting.\\
\textbf{Ukraine / EU (discussion frame):} accession blocked/conditional on \textbf{security}/\textbf{war}, macro \textbf{instability}, and \textbf{rule of law} (corruption/\textbf{judicial} reform)---parallel \textbf{Article 2} values / \textbf{governance} screens in EU courses.\\
\textbf{Closing moral (lecture):} \textbf{technology} and \textbf{capital} necessary, but \textbf{human motivation} is the \textbf{engine} through the ``\textbf{nightmare}'' of \textbf{cultural adaptation} and long \textbf{commitment}. \textbf{Ask questions} early: youth is not knowing everything. \textbf{Silence} buys \textbf{costly mistakes}.\\
\textbf{Exam link:} label \textbf{strategy type}; sequence \textbf{integration} stages; define \textbf{static} vs \textbf{dynamic} effects; tie \textbf{ECI} to \textbf{FDI} site choice + \textbf{people} readiness.
}

\cornell{19) Integration Drivers, Hierarchy (with PTA), Location Screen, and Execution Risk (Lecture Thread)}
{\textbf{Strategy horizon:} \textbf{15 days} vs \textbf{15 years}?\\
\textbf{Integration:} who wins/loses inside the bloc?\\
\textbf{Trade creation} vs \textbf{diversion}---plain language?\\
\textbf{Hierarchy:} step before \textbf{FTA}?\\
\textbf{Location:} what screens beyond \textbf{ECI}?\\
\textbf{Phase 3 execution:} where do \textbf{technology} and \textbf{capital} leak?\\
\textit{Realist close:} where do \textbf{ports}/\textbf{infrastructure} make the business?}
{
\textbf{Time discipline (aphorism):} \textbf{strategy} framed as \textbf{long-run} (\textbf{15-year} consistency) not \textbf{short} \textbf{tactical} sprints (\textbf{15-day} mindset); patience and steady execution beat reset churn (echo blocks~10, 17).\\
\textbf{Economic integration (political economy):} \textbf{two or more} states agree to \textbf{reduce barriers}; driven by both \textbf{politics} and \textbf{economics}; \textbf{prosperity} aim for members but \textbf{gains need not be equal}---distributional fights inside blocs matter for firm and voter risk.\\
\textbf{Static trade effects (lecture phrasing):} \textbf{trade creation}---\textbf{efficient} exchange \textbf{among} members as barriers fall; \textbf{trade diversion}---imports \textbf{swing} from a \textbf{lower-cost outsider} to a \textbf{higher-cost insider} because the union discriminates---welfare ambiguous, politics hot (block~18 vocabulary).\\
\textbf{Hierarchy (full rung list):} (0) \textbf{preferential trading agreement (PTA)}---narrow, selective concessions (\textbf{simplest} cooperation step); then \textbf{FTA} $\rightarrow$ \textbf{customs union} $\rightarrow$ \textbf{common market} (factors + \textbf{services} mobility in lecture list) $\rightarrow$ \textbf{economic \& monetary union} (\textbf{shared currency}/monetary policy, \textbf{EU} anchor).\\
\textbf{Location choice bundle:} \textbf{corporate} ``where'' hinges not only \textbf{ECI}/\textbf{product--competency fit} but \textbf{infrastructure} (\textbf{ports}, corridors---\textbf{Mexico}/\textbf{Chile} color), \textbf{political stability}, and \textbf{regulatory safety}/predictability; mis-match any leg and \textbf{execution} cost spikes.\\
\textbf{FDI reminder:} physical \textbf{capital} and \textbf{operations} move---\textbf{people} and \textbf{departments}, not only boxes shipped; higher \textbf{complexity} than \textbf{exporting} (blocks~13, 15).\\
\textbf{Human capital:} no \textbf{internationalization} without \textbf{internationalized people}; \textbf{expats} protect \textbf{technology} and \textbf{capital} but cost \textbf{cash}---trust and \textbf{governance} on site non-negotiable.\\
\textbf{Execution phase (block~10 echo):} phases~1--2 feel \textbf{theoretical}; \textbf{phase~3} is where \textbf{cash burns} and \textbf{risk peaks}---firms lose \textbf{IP}/\textbf{know-how} and \textbf{money} if local control and trusted leads fail; ``\textbf{realist}'' view---go where there is \textbf{real demand} and the \textbf{infrastructure} fits your \textbf{product}, not where a map looks pretty.\\
\textbf{Anecdote thread:} instructor ties choices to operating memory in \textbf{Mexico} and \textbf{Chile}---use as illustration of \textbf{port}/\textbf{logistics} reality checks, not universal law.\\
\textbf{Exam link:} order \textbf{PTA}$\rightarrow$\textbf{EMU}; define \textbf{creation}/\textbf{diversion}; tie \textbf{location screen} to \textbf{phase~3} risk and staffing.
}

\cornell{20) International PM, Commitment, Authority, and U-Curve Readiness (Lecture Thread)}
{\textbf{Plan} vs \textbf{finish}---what does the professor privilege?\\
\textbf{Commitment} vs \textbf{skill} for int'l seats?\\
\textbf{Phase 3}---where is \textbf{capital} most exposed?\\
\textbf{Job rotation}---why before \textbf{expat} leadership?\\
\textbf{Authority} beyond the \textbf{title}?\\
\textit{Failure drivers:} bad \textbf{site} vs untrusted \textbf{people}?}
{
\textbf{Execution thesis:} \textbf{strategy} is worthless without \textbf{delivery}---the job is to \textbf{finish}; many int'l programs fail from \textbf{inconsistency} and from never truly \textbf{securing} the early investment through \textbf{phase~3} (blocks~10--12).\\
\textbf{Commitment $>$ raw skill (lecture):} \textbf{engineering}/\textbf{finance} matter, but \textbf{commitment}/\textbf{loyalty} to \textbf{HQ} and the \textbf{mission} outrank credentials alone for protecting \textbf{capital} abroad; ``\textbf{trusted}'' \textbf{HQ} staff often preferred initially to pure locals because alignment with parent goals is clearer (trade-off: \textbf{cost}, speed to context).\\
\textbf{Financial discipline:} manager as \textbf{steward} of \textbf{capital}---\textbf{cash generation} via \textbf{cost} control and \textbf{revenue} lift while holding \textbf{margin} (e.g.\ \textbf{20\%} reference rows from earlier sessions---map to your case numbers).\\
\textbf{Job rotation (HR prep):} cycle leads through \textbf{production}, \textbf{sales}, \textbf{HR}/staff functions so international \textbf{GM} sees the \textbf{whole firm}---not silo \textbf{optimization} only (vocabulary link).\\
\textbf{Authority as respected status:} especially in high-\textbf{PDI} markets (\textbf{Mexico}, \textbf{Chile} in stories), a \textbf{young} manager may lack automatic \textbf{deference}; \textbf{authority} is \textbf{earned} via \textbf{performance}, relationship care inside the \textbf{chain of command}, and \textbf{CQ}---title alone does not close deals (block~16 echo).\\
\textbf{IHRM bundle:} \textbf{expatriates} carry \textbf{company DNA} and guard \textbf{technology} transfer; expensive package demands \textbf{U-curve} literacy---\textbf{honeymoon} $\rightarrow$ \textbf{shock} (language, \textbf{logistics}, expectation gaps) $\rightarrow$ \textbf{adjustment} (local rules, \textbf{empathy}) $\rightarrow$ \textbf{mastery}/steady state; early quit often tracks \textbf{shock} mismanagement (block~9).\\
\textbf{Lifecycle recap:} \textbf{Phases 1--2}---research, \textbf{location}, \textbf{market} analysis (\textbf{planning}); \textbf{Phase~3}---\textbf{resources} move, \textbf{hiring}, \textbf{production} start (\textbf{budget} danger); \textbf{Phase~4}---\textbf{profitability}, \textbf{relationship} depth (block~10 numbering).\\
\textbf{Anecdote frame:} instructor contrasts \textbf{France} classroom logic with \textbf{Mexico}/\textbf{Chile} \textbf{ground truth}---theory meets \textbf{operational friction}.\\
\textbf{Why int'l efforts fail (lecture):} often \textbf{not} a bad \textbf{product} but \textbf{poor location choice} and/or \textbf{lack of trusted on-site management}; cannot \textbf{copy-paste} a \textbf{French} playbook into \textbf{Chile}---need \textbf{physical presence}, \textbf{cultural adaptation}, and a lead \textbf{fully aligned} with \textbf{HQ} to \textbf{shield} \textbf{capital}.\\
\textbf{Exam link:} chain \textbf{SWOT}/\textbf{location} $\rightarrow$ \textbf{phase plan} $\rightarrow$ \textbf{people} (\textbf{commitment}, \textbf{rotation}, \textbf{expat}) $\rightarrow$ \textbf{U-curve} risk controls.
}

\section*{Summary Block (Cornell Bottom)}
\begin{itemize}
  \item This course needs highest practical focus because it sits at AP threshold.
  \item Write answers with clear links: context -> mode choice -> execution risk -> mitigation.
  \item Aim to move from descriptive to analytical responses with explicit trade-offs.
  \item Chain + culture: pair Porter activity choices with Hofstede/HR risk; avoid home-country templates abroad.
  \item Entry: \textbf{B vs C}---empty markets often hide missing demand/infrastructure; GDP and political risk screen before mode choice.
  \item Cross-culture: \textbf{PDI/CQ}, U-curve, and non-verbal communication---link behavior to dimension, not stereotypes only.
  \item Phases 1--2 before execution: \textbf{artificial} (systems) vs \textbf{natural} (culture/lang) risk; $\sim$6-year maturity, 2-year execution cash burn.
  \item PM/HR: \textbf{expat vs local} trade-offs; \textbf{margin} green/red and portfolio average; \textbf{commitment}/trust with capital abroad.
  \item Cash/CRM: \textbf{50\%} down-payment discipline; yellow/red margin bands; retention; \textbf{6-month} political risk vs long-term plan.
  \item Economic complexity: \textbf{ECI}-style screening, stuck-in-middle trap, diversification vs specialization; location--competency fit; matrix (vertical + horizontal).
  \item Governance/strategy: corporate--competitive--functional pyramid, \textbf{SWOT} audit, \textbf{IP}/oversight; realist triad (power, financial control, resilience).
  \item Int'l spectrum: motives (market / resource / efficiency); domestic $\rightarrow$ export $\rightarrow$ subsidiary $\rightarrow$ integration; direct vs indirect export; \textbf{2--3}-year cash shield.
  \item IHRM: global vs multi-domestic; commitment/trust vs pedigree; \textbf{Maslow}/\textbf{Vroom}; team diversity trade-offs; negotiator \textbf{authority}; \textbf{CQ} vs pay in hard markets.
  \item Applied studio (\textbf{Caf\'e Ol\'e}): \textbf{SWOT}+\textbf{PESTEL}, student pivot pricing/membership, contingencies, Q1--Q4 rollout; \textbf{slow is fast} (\textbf{15} years not \textbf{15} days); align wording with instructor.
  \item Strategy/integration map: \textbf{transnational} hybrid; FTA $\rightarrow$ customs union $\rightarrow$ common market $\rightarrow$ \textbf{EMU}; static (\textbf{creation}/\textbf{diversion}) vs dynamic effects; \textbf{ECI}/competencies; \textbf{triad}/\textbf{G20}; \textbf{Ukraine} governance barrier; motivation + ask questions.
  \item Integration + entry (block 19): \textbf{PTA} before FTA; politics/economics; unequal gains; location = \textbf{ECI} + \textbf{infrastructure} + stability + \textbf{regulatory safety}; phase 3 \textbf{execution} = peak leak of tech/capital; realist fit to \textbf{ports}/demand.
  \item IPM/IHRM execution (block 20): plan vs \textbf{finish}; \textbf{commitment} vs skill; \textbf{capital}/\textbf{margin}; \textbf{job rotation}; \textbf{authority} as earned status; \textbf{U-curve}; failure = bad \textbf{location} or untrusted \textbf{management}; no \textbf{home-country copy-paste} (\textbf{France} $\neq$ \textbf{Chile}).
  \item Finals tile status: no course-specific image tiles detected in current upload batch; keep updating from class notes and slide packets.
\end{itemize}

\end{document}
